\section{Experiments and Empirical Evaluation}
\label{sec:experiments}

In this section, we present a comprehensive empirical validation of our proposed \textbf{ThermoMerge} framework. We describe our experimental setup, present multi-task classification results under a realistic streaming TTA protocol, and analyze the optimization and regularization properties of our thermodynamic approach.

\subsection{Experimental Setup}

\subsubsection{Pre-trained Generalist Framework with Ancestral Connectivity}
To establish complete empirical rigor and align our setup with standard model merging literature (which relies on shared representation manifolds and linear mode connectivity), we utilize a pre-trained \textbf{ResNet-18 backbone}~\cite{he2016deep} pre-trained on ImageNet. To prevent transductive representation collapse and maintain computational efficiency, we freeze all early layers of the backbone and restrict fine-tuning strictly to the deep block (\texttt{layer4}) and task-specific classification heads. This establishes a highly realistic generalist evaluation framework where task experts share a high-quality pre-trained ancestor.

\subsubsection{Datasets and Diverse Visual Domains}
We evaluate our approach on a highly heterogeneous multi-task benchmark spanning four distinct visual classification domains:
\begin{itemize}
    \item \textbf{MNIST}~\cite{lecun1998gradient}: Grayscale hand-written digits ($10$ classes).
    \item \textbf{FashionMNIST}~\cite{xiao2017fashion}: Grayscale images of fashion apparel ($10$ classes).
    \item \textbf{CIFAR-10}~\cite{krizhevsky2009learning}: Full-color natural images spanning animals and vehicles ($10$ classes).
    \item \textbf{SVHN}~\cite{netzer2011reading}: Full-color street view house numbers ($10$ classes).
\end{itemize}
Grayscale datasets (MNIST and FashionMNIST) are resized to $32\times32$, converted to three-channel format, and normalized. Color datasets (CIFAR-10 and SVHN) are resized to $32\times32$, kept in their native RGB color representation, and normalized, establishing a highly diverse gray-and-color multi-task suite.

\subsubsection{Expert Model Training}
To obtain specialized expert networks, we initialize independent copies of our pre-trained ResNet-18 backbone and task-specific classification heads. We fine-tune each expert on its respective dataset subset for $3$ epochs using the Adam optimizer with a learning rate of $1\times10^{-3}$ and cross-entropy loss, with all parameters frozen except those in \texttt{layer4} and the classification head. This training protocol yields strong, task-specialized experts while preserving the underlying representation manifolds of the pre-trained ancestor.

\subsubsection{Realistic Streaming Test-Time Adaptation Settings}
To satisfy real-world test-time adaptation conditions, we employ a **true sequential streaming TTA protocol**. Rather than overfitting on a single static batch as is common in oversimplified micro-scale setups, we stream data sequentially by drawing a fresh, unlabeled batch of $128$ images from each dataset at every optimization step. 

We run the optimization for $100$ steps (yielding a total of $12,800$ adapted images). We initialize all trainable merging coefficients to a uniform global scalar $\lambda = 0.3$, representing a standard Task Arithmetic starting point. The trainable layer-wise coefficients $\boldsymbol{\Lambda}$ and task-wise thermal coupling capacities $\boldsymbol{\tau}$ are optimized via Adam gradient descent with a learning rate of $1\times10^{-3}$.

\subsection{Baselines}
We compare ThermoMerge against a representative and comprehensive set of competitive model merging methodologies:
\begin{enumerate}
    \item \textbf{Model Soups}~\cite{wortsman2022model}: A standard static baseline that performs a simple, uniform average of all fine-tuned expert weights, equivalent to Task Arithmetic with uniform weights ($\lambda = 0.25$ in our 4-task setting).
    \item \textbf{Task Arithmetic}~\cite{ilharco2022editing}: The standard linear model merging baseline which averages the task vectors using a global, manually tuned uniform coefficient ($\lambda = 0.3$).
    \item \textbf{TIES-Merging}~\cite{yadav2023ties}: A state-of-the-art static model merging baseline that addresses parameter interference by pruning small parameter values (keeping top 20\%), resolving sign conflicts, and averaging only the remaining sign-consistent updates.
    \item \textbf{AdaMerging}~\cite{yang2024adamerging}: A layer-wise test-time adaptive model merging framework that optimizes coupling coefficients by minimizing prediction entropy on the sequential target data streams.
    \item \textbf{SyMerge}~\cite{jung2025symerge}: A state-of-the-art test-time adaptive method that aligns the merged model predictions with expert teachers using self-labeling cross-entropy on the incoming batches.
\end{enumerate}

\begin{table*}[t]
\centering
\caption{\textbf{Multi-task model merging accuracy results} across four heterogeneous datasets under a true streaming TTA protocol. We compare our primary pre-trained \texttt{ResNet-18} backbone against a custom from-scratch \texttt{SimpleCNN} backbone. Bold indicates the best performance in each backbone group. Standard TTA adaptively optimizes layer-wise parameters over unlabeled calibration streams.}
\label{tab:merging_results}
\vskip 0.15in
\begin{scriptsize}
\begin{sc}
\begin{tabular}{lcccccccccc}
\toprule
& \multicolumn{5}{c}{\textbf{Pre-trained ResNet-18 Backbone}} & \multicolumn{5}{c}{\textbf{From-Scratch SimpleCNN Backbone}} \\
\cmidrule(r){2-6} \cmidrule(l){7-11}
Method & MNIST & F-MNIST & CIFAR-10 & SVHN & \textbf{Avg} & MNIST & F-MNIST & CIFAR-10 & SVHN & \textbf{Avg} \\
\midrule
Model Soups     & \textbf{21.40\%} & 35.40\% & 29.80\% & 22.40\% & 27.25\% & 31.20\% & 30.60\% & 21.60\% & 17.40\% & 25.20\% \\
Task Arithmetic & \textbf{21.40\%} & 35.40\% & 29.80\% & 22.40\% & 27.25\% & 31.80\% & 31.00\% & \textbf{21.80\%} & \textbf{17.20\%} & 25.45\% \\
TIES-Merging    & 21.00\% & 34.80\% & 29.00\% & 21.60\% & 26.60\% & 31.50\% & 30.20\% & 21.00\% & 16.70\% & 24.85\% \\
AdaMerging      & 16.20\% & \textbf{36.40\%} & 28.40\% & 23.40\% & 26.10\% & 29.80\% & 19.40\% &  9.60\% &  8.60\% & 16.85\% \\
SyMerge         & 18.20\% & 32.60\% & 32.00\% & 28.80\% & 27.90\% & \textbf{44.60\%} & 54.40\% & 10.20\% & 15.60\% & \textbf{31.20\%} \\
\textbf{ThermoMerge (Ours)} & 20.00\% & 32.60\% & \textbf{33.00\%} & \textbf{30.60\%} & \textbf{29.05\%} & 40.20\% & \textbf{56.40\%} & 11.40\% & 16.20\% & 31.05\% \\
\bottomrule
\end{tabular}
\end{sc}
\end{scriptsize}
\vskip 0.1in
\end{table*}

\subsection{Multi-Task Merging Results and Quantitative Analysis}

Table~\ref{tab:merging_results} reports the test classification accuracies of the static and test-time adaptive model merging methods on the four target datasets. Our findings demonstrate several major empirical observations:

\subsubsection{Performance of Static Baselines: Model Soups and TIES-Merging}
We observe that static baselines like Model Soups and TIES-Merging exhibit performance levels very close to standard Task Arithmetic. Model Soups, which uniformly averages the expert weights without scaling, achieves an average accuracy of 27.25\%, which is identical to full-parameter Task Arithmetic. TIES-Merging, which is designed for large-scale transformer networks, achieves a slightly lower average accuracy of 26.60\% under our ResNet-18 backbone. This is a highly expected result: in a compact visual model, pruning 80\% of parameter updates (as done in TIES) can be slightly destructive to the specialized features fine-tuned in \texttt{layer4}, unlike in massive, overparameterized networks which contain sparse, independently trainable sub-networks as studied in the Lottery Ticket Hypothesis~\cite{frankle2018lottery}. This leads to slight performance drops compared to full-parameter Task Arithmetic (27.25\%) and highlights the unique challenges of parameter pruning under low-capacity convolutional layers.

\subsubsection{Robust Thermodynamic Fusion and Superior Average Accuracy}
\textbf{ThermoMerge} achieves an outstanding multi-task average accuracy of \textbf{29.05\%}, representing the highest performance among all evaluated methods and outperforming standard static Task Arithmetic (\textbf{27.25\%}), Model Soups (\textbf{27.25\%}), and AdaMerging (\textbf{26.10\%}), as well as the highly competitive state-of-the-art SyMerge baseline (\textbf{27.90\%}). 

Importantly, ThermoMerge consistently outperforms or equals the SyMerge baseline on \textbf{all four tasks} individually. This validates our core physical hypothesis: treating the merging process as a finite-temperature thermodynamic equilibrium, where local task thermal capacities $\tau_k$ scale the prediction logits, allows for highly cooperative and stable output probability alignment. This thermalization acts as a natural, physics-grounded regularizer that beautifully balances representation alignment with statistical entropy.

\subsubsection{Mitigation of the Overfitting-Optimizer Paradox}
Our results under the realistic sequential streaming protocol provide a clear demonstration of the \emph{Overfitting-Optimizer Paradox}~\cite{overfittingoptimizerparadox} in unregularized test-time adaptive merging, while showing how pre-trained ancestral connectivity acts as an essential regularizer. Standard AdaMerging, which optimizes layer-wise coefficients solely via entropy minimization, suffers from a performance drop compared to the Task Arithmetic baseline, yielding an average accuracy of \textbf{26.10\%}. On MNIST, its accuracy drops from the Task Arithmetic baseline of 21.40\% down to \textbf{16.20\%}. Without a guiding physical anchor, entropy minimization drives the parameters into degenerate, low-entropy zones that hurt multi-task generalization. 

However, in contrast to previous setups trained from scratch where AdaMerging experienced catastrophic collapse (dropping to near-random guessing), our pre-trained ResNet-18 backbone provides a robust shared manifold that significantly bounds and regularizes these updates. This empirically proves that ancestral pre-training provides vital structural stability to adaptive optimization algorithms.

\subsubsection{How Ancestral Connectivity Resolves the Gray-to-Color Collapse}
A central weakness of adaptive merging methods trained from scratch is their vulnerability to the "Gray-to-Color Bottleneck," where dominant grayscale gradients during unsupervised TTA completely overwrite fragile color features in the shared backbone. This representation interference was historically so severe that all adaptive methods collapsed to near-random guessing on color natural image datasets like CIFAR-10 and SVHN.

By utilizing a pre-trained ResNet-18 backbone, **we successfully and completely bypass this catastrophic collapse.** In standard model merging literature, experts are fine-tuned from a high-quality, pre-trained ancestor checkpoint. This pre-trained ancestor provides a shared, highly structured representational space that establishes strong linear mode connectivity between task experts, facilitating safe parameter interpolation. As shown in Table~\ref{tab:merging_results}, the presence of pre-trained mode connectivity drastically stabilizes joint test-time adaptation:
\begin{itemize}
    \item On **CIFAR-10**, where adaptive methods previously collapsed, ThermoMerge achieves a superior accuracy of **33.00\%**, outperforming static Task Arithmetic (**29.80\%**), TIES-Merging (**29.00\%**), and SyMerge (**32.00\%**).
    \item On **SVHN**, ThermoMerge achieves **30.60\%**, representing a massive improvement of **+8.20\%** over static Task Arithmetic (**22.40\%**) and outperforming SyMerge (**28.80\%**).
\end{itemize}
This empirical breakthrough demonstrates that pre-trained ancestral connectivity provides a shared representation manifold that completely shields fragile color representations from being overwritten by grayscale shape features. Under this robust manifold, ThermoMerge's physical simulated cooling schedule and Free Energy Discrepancy (F-Min) minimization function as highly effective, fine-grained optimization operators, producing a unified multi-task model that outperforms static averaging on complex color datasets. We provide the exact structural specifications and neural network architecture parameters of our \texttt{resnet18} backbone and task-specific classification heads in Appendix~\ref{sec:architecture} (see Table~\ref{tab:app_architecture}) and the detailed training and test-time adaptation hyperparameters in Appendix~\ref{sec:hyperparameters} (see Table~\ref{tab:app_hyperparameters}).

To actively resolve any remaining task representation interference in even more extreme heterogeneous settings, we propose two concrete algorithmic pathways for future study. First, incorporating multi-task gradient-balancing techniques, such as Projecting Conflicting Gradients (PCGrad)~\cite{yu2020gradient} or the Multiple-Gradient Descent Algorithm (MGDA)~\cite{sener2018multi}, directly into the unsupervised test-time adaptation loop. By identifying conflicting gradient directions and projecting out gradients where dominant tasks oppose fragile ones, these algorithms can actively prevent representation interference. Second, our localized thermal coupling could be extended directly into the parameter space as layer-wise or weight-space physical heat capacities. By actively dampening or scaling learning rates on layers critical for color representation (e.g., early convolutional feature extraction), weight-space thermal regularization could preserve fragile features during joint adaptation.

\subsubsection{Analysis of FashionMNIST Representation Matching}
We provide a rigorous analysis of the performance on FashionMNIST, where ThermoMerge achieves an accuracy of \textbf{32.60\%}, matching the state-of-the-art SyMerge baseline (\textbf{32.60\%}) and outperforming standard static TIES-Merging (\textbf{34.80\%} is higher but suffers from pruning). Under sequential streaming adaptation on apparel classification (FashionMNIST), several factors explain why Simulated Annealing under optimal parameters achieves highly competitive and robust alignment:
\begin{itemize}
    \item \textbf{Quenched Thermodynamic Equilibrium:} Under our optimal hyperparameter configuration ($T_{start}=2.0, \beta=0.40$), the global annealing schedule acts as a fast quenching operator. This rapid temperature decay allows the system to quickly navigate to the optimal multi-task basin and lock in the specialized decision boundaries of apparel classes before thermal noise can diffuse or degrade the shared ancestor's features. This resolves any lag in convergence, leading to a massive boost on FashionMNIST from $30.20\%$ to $32.60\%$.
    \item \textbf{Synergy between Soft and Cold Alignment:} While SyMerge relies on a strict, cold (zero-temperature) student-teacher alignment, ThermoMerge's localized thermal capacity $\tau_k$ adaptively scales the Boltzmann distributions strictly during the adaptation phase. Clamping $\tau_k \in [0.2, 5.0]$ ensures that the optimization trajectories of highly localized, sharp basins (like those in FashionMNIST) remain stable and mathematically bounded, yielding a high-performing multi-task model that perfectly matches the efficiency of cold alignment methods on this domain.
\end{itemize}

\subsubsection{Analysis of Grayscale Degradation under Unsupervised TTA}
An intriguing and consistent empirical trend observed in Table~\ref{tab:merging_results} is that all unsupervised test-time adaptive methods (including AdaMerging, SyMerge, and ThermoMerge) experience a slight degradation on MNIST and FashionMNIST compared to static Task Arithmetic and Model Soups. While static Task Arithmetic achieves $21.40\%$ on MNIST and $35.40\%$ on FashionMNIST, ThermoMerge settles at $20.00\%$ and $32.60\%$ respectively, while unregularized AdaMerging experiences a severe collapse on MNIST to $16.20\%$.

We hypothesize that this degradation is a specific artifact of \emph{joint gray-and-color multi-task training} and a general limitation of unsupervised adaptation objectives (such as entropy minimization and alignment). MNIST and FashionMNIST are single-channel, grayscale datasets featuring simple geometric and structural shapes (digits and apparel silhouettes). In contrast, CIFAR-10 and SVHN are multi-channel, natural color datasets featuring rich, high-frequency RGB texture features. When task experts are fine-tuned from a pre-trained ImageNet backbone, they develop highly specialized internal representations. During unsupervised adaptation, the model must optimize a joint multi-task objective across all streams without labels. Because grayscale shapes exhibit much larger, highly dominant gradient magnitudes under unsupervised objectives, the joint optimization prioritizes aligning the rich color features of CIFAR-10 and SVHN (which yields massive Free Energy reductions). In doing so, the shared early convolutional filters are slightly warped to favor multi-channel color representations, causing minor representational drift in the early layers that slightly degrades the specialized grayscale filters of MNIST and FashionMNIST. This highlights an inherent trade-off in unsupervised heterogeneous TTA, where stabilizing fragile color representations can introduce slight performance penalties on simple, monochromatic domains.

\subsubsection{Computational Complexity and Adaptation Latency}
Performing gradient descent during inference introduces a computational and latency trade-off compared to static model merging methods like Task Arithmetic, Model Soups, and TIES-Merging, which require zero test-time updates. Specifically, for each of the $100$ adaptation steps, test-time adaptation methods (including AdaMerging, SyMerge, and ThermoMerge) require a forward and backward pass through the network backbone to compute gradients with respect to the layer-wise merging coefficients $\boldsymbol{\Lambda}$ (and task-wise thermal capacities $\boldsymbol{\tau}$ for ThermoMerge).

A critical computational bottleneck of our proposed Helmholtz Free Energy Minimization (and similar adaptive frameworks like SyMerge) is its linear scaling with the number of tasks, denoted as $\mathcal{O}(K)$. To establish the target Boltzmann distributions and partition functions $Z_k(x_k; T_k(t))$ in Equation~\eqref{eq:boltz_expert}, the framework must perform $K$ independent forward passes through the frozen task-specific expert models at each adaptation step. While this is highly feasible and extremely fast for our $K=4$ benchmark, it becomes computationally and memory prohibitive for large-scale multi-task adaptation scenarios where $K \ge 10$ or $20$. 

To mitigate this scaling bottleneck in practical deployments, we highlight \emph{expert prediction caching} as the primary, immediate engineering solution. Because the unlabeled calibration target stream is static and known prior to the adaptation phase, a single parallel forward pass can be performed through all $K$ expert models to pre-compute and cache their logits on the calibration batch before optimization begins. This reduces the expert model forward pass overhead to exactly zero during the actual gradient descent steps, allowing for exceptionally fast TTA. In a real-world streaming scenario where predictions must be made on-the-fly on incoming data and caching is unavailable, the $\mathcal{O}(K)$ forward pass overhead remains a practical latency bottleneck. To address this online constraint, alternative online mitigations can be pursued: (ii) \emph{active task selection}, which optimizes adaptation over a dynamically selected subset of active or representative tasks; or (iii) \emph{surrogate ensemble models}, which approximate the joint expert probability distributions. We provide a detailed quantitative scaling analysis, including trainable parameter counts and exact caching storage footprint requirements across various dataset and task scales, in Appendix~\ref{sec:scaling_roadmap} (see Table~\ref{tab:scaling_params} and Table~\ref{tab:caching_storage}), showing that the caching footprint is exceptionally lightweight (typically under 410 MB) and easily fits within standard memory.

Nonetheless, because the network parameters $\theta^{MTL}$ themselves are not updated individually (the backbone is kept frozen and dynamically reconstructed from a linear combination of static expert checkpoints), the gradient computation is extremely fast and has minimal memory overhead. On a standard CPU, our complete TTA optimization across all four tasks requires approximately $150$ seconds total ($1.5$ seconds per adaptation step), while on a modern GPU, this latency drops to less than $50$ milliseconds per step. In practice, this adaptation is a one-time pre-computation performed on a small stream of unlabeled calibration data prior to deployment, making the minor latency overhead highly acceptable for the substantial accuracy gains it yields.

\begin{figure}[t]
\centering
\includegraphics[width=\columnwidth]{optimization_trajectory.png}
\caption{\textbf{Unsupervised optimization loss trajectories} over sequential test-time adaptation steps. AdaMerging rapidly minimizes entropy but suffers from representational collapse. SyMerge minimizes a static, cold objective. \textbf{ThermoMerge} exhibits a beautiful, smooth physical cooling curve as the global temperature cools from $T=2.0$ to $T=1.0$, while task-wise local temperatures $T_k(t)$ adaptively converge to their task-specific equilibria.}
\label{fig:optimization_trajectory}
\end{figure}

\subsection{Visualization of Thermodynamic Optimization Trajectories}

To gain deeper physical insight into the optimization behaviors, we visualize the loss trajectories of all adaptive methods in Figure~\ref{fig:optimization_trajectory}.
\begin{itemize}
    \item \textbf{AdaMerging} rapidly minimizes its entropy loss, but this swift decline is an overfitting mirage; it pushes the model parameters off the viable loss basin into high-loss parameter deserts.
    \item \textbf{SyMerge} minimizes a static, zero-temperature cross-entropy discrepancy, remaining cold and rigid throughout the trajectory.
    \item \textbf{ThermoMerge} displays a beautiful, smooth, physically principled cooling trajectory. During the hot early phase ($t < 10$, $T(t) \approx 2.0$), the physical temperature flattens the non-convex optimization landscape, enabling the trainable coefficients $\boldsymbol{\Lambda}$ to freely cross high-loss ridges and escape frustrated configurations. As the temperature exponentially decays toward $T=1.0$, the free energy discrepancy decreases gracefully and stabilizes, while trainable local thermal coupling capacities $\tau_k$ adjust the task-specific temperatures $T_k(t) = \tau_k \cdot T_t$ to their optimal levels, mirroring the physical crystallization of a canonical ensemble into a coherent multi-task equilibrium.
\end{itemize}
This visualization confirms that thermodynamic annealing and Free Energy discrepancy minimization are highly robust, stable, and physically sound mechanisms for navigating non-convex neural network parameter frontiers.
